% ch16-systems.tex -- Vectorisation and data parallelism
\chapter{Vectorisation and data parallelism}
\label{ch:systems16}
\index{SIMD}\index{NEON}\index{vectorisation}

Threading splits work across cores; vectorisation splits it across
lanes in one instruction. This is the third parallelism the book
has built, and the cheapest: change nothing about the algorithm,
get 4 items per issue. The catch is in the chapter's second
measurement --- the compiler may already have taken the win.

\section{Data-level parallelism}
\begin{figure}[htbp]
\centering
\begin{tabular}{lll}
\toprule
parallelism & unit & book chapter \\
\midrule
instruction-level (ILP) & instrs/cycle in one thread & ch6--7 \\
data-level (DLP) & lanes/instr --- this chapter & ch16 \\
thread-level (TLP) & cores & ch13 \\
\bottomrule
\end{tabular}
\caption{The three parallelisms the book has now covered. DLP is
free ILP: the pipeline issues once and 4 lanes retire --- the
widest speedup-per-line-changed in systems work.}
\label{fig:dlp16}
\end{figure}


\section{One instruction, four lanes}
\begin{figure}[htbp]
\centering
\begin{tabular}{lccc}
\toprule
& \multicolumn{3}{c}{one instruction, four lanes} \\
\midrule
scalar & y0+ax0 & y1+ax1 & y2+ax2 \dots \\
NEON   & \multicolumn{3}{c}{\texttt{vmlaq}: lanes 0--3 in one op} \\
\bottomrule
\end{tabular}
\caption{Data-level parallelism: one \texttt{vmlaq} updates 4 floats
per issue where the scalar loop updates one. The ALU count is the
same; the instruction stream is $4\times$ shorter.}
\label{fig:lanes16}
\end{figure}


\section{SAXPY, measured}
\begin{figure}[htbp]
\centering
\begin{tabular}{lrrr}
\toprule
n & scalar (ms) & NEON (ms) & speedup \\
\midrule
$2^{16}$ & 0.025 & 0.007 & 3.5$\times$ \\
$2^{19}$ & 0.189 & 0.064 & 3.0$\times$ \\
$2^{22}$ & 1.384 & 0.331 & 4.2$\times$ \\
$2^{25}$ & 8.520 & 2.885 & 3.0$\times$ \\
\bottomrule
\end{tabular}
\caption{Real \texttt{ch16\_saxpy} sweep, \texttt{-fno-vectorize}
scalar vs NEON intrinsics: 3--4$\times$ at large n --- close to the
4-lane ideal, capped by memory bandwidth at the biggest sizes.}
\label{fig:saxpy16}
\end{figure}


\begin{example}[3--4$\times$ real]
\texttt{-fno-vectorize} scalar vs NEON intrinsics across
$n=2^{10}\dots2^{25}$: speedup tracks the 4-lane ideal (3.5$\times$
at $2^{16}$, 4.2$\times$ at $2^{22}$), sagging to 3.0$\times$ at
$2^{25}$ where memory bandwidth, not issue width, is the limit ---
ch9's roofline arriving on schedule.
\end{example}

\section{The folklore item that fails}
\begin{figure}[htbp]
\centering
\begin{tabular}{lrr}
\toprule
variant & ms at $n=2^{25}$ & verdict \\
\midrule
scalar \texttt{-fno-vectorize} & 8.520 & real scalar \\
scalar \texttt{-O3} (auto-vec) & 2.891 & compiler SIMD'd it \\
NEON intrinsics & 2.885 & same \\
\bottomrule
\end{tabular}
\caption{The chapter's folklore entry dies here: ``intrinsics beat
the compiler'' measured $1.00\times$ at $n=2^{25}$ --- clang's
auto-vectoriser emits equivalent NEON. Write intrinsics for
control, not for speed the compiler already finds.}
\label{fig:autovec16}
\end{figure}


\begin{example}[Intrinsics $=$ compiler, measured]
At \texttt{-O3} clang auto-vectorises the scalar SAXPY: 2.891\,ms
vs the hand-written intrinsics' 2.885\,ms --- $1.00\times$. The
``intrinsics beat the compiler'' claim dies on this machine.
Write intrinsics for guaranteed vectorisation on compilers that
won't, or for ops the compiler can't express --- not for raw speed
it already finds.
\end{example}

\section{SoA vs AoS: the layout is the vector}
\begin{figure}[htbp]
\centering
\begin{tabular}{lrr}
\toprule
layout & x+=m over 16M particles & memory touched \\
\midrule
AoS (struct\,\{x,y,z,m\}) & 4.7\,ms & 4 arrays' worth \\
SoA (x[], y[], z[], m[]) & 2.0\,ms & 2 arrays only \\
\bottomrule
\end{tabular}
\caption{Real \texttt{ch16\_soa}: the same update is 2.35$\times$
faster as SoA because a cache line fetches only the two fields used.
AoS drags y and z through the hierarchy for nothing.}
\label{fig:soa16}
\end{figure}


A struct-of-arrays update reads only the fields used; array-of-
structs drags the whole record through every cache line. Measured:
2.0\,ms vs 4.7\,ms on 16M particles --- $2.35\times$ for a pure
layout change, no SIMD written. The vector unit loves SoA because
the memory system already does.

\section{The three vector ISAs}
\begin{figure}[htbp]
\centering
\begin{tabular}{ll}
\toprule
ISA & model \\
\midrule
NEON (arm64) & 128-bit fixed lanes \\
AVX-512 & 512-bit + masks \\
RVV (RISC-V) & vector-length-agnostic, scalable \\
\bottomrule
\end{tabular}
\caption{Three SIMD philosophies. RVV's trick is portability across
hardware widths: the same loop runs on 128-bit and 1024-bit
implementations --- NEON's fixed 128 is the price of its simplicity.}
\label{fig:rvv16}
\end{figure}


\section{Common mistakes}
\begin{itemize}
\item \textbf{Assuming the scalar baseline is scalar.} At
      \texttt{-O3} the compiler vectorises anyway --- our honest
      baseline needed \texttt{-fno-vectorize} to be honest.
\item \textbf{``Intrinsics always win.''} $1.00\times$ vs
      auto-vec here; the hand-written version earns its keep in
      control and portability, not raw speed.
\item \textbf{Vectorising the wrong loop.} Amdahl again --- a
      4$\times$ kernel at 10\% of runtime buys 3\%.
\item \textbf{AoS by default.} The struct grouping that's natural
      to write is the layout that costs $2.35\times$ ---
      SoA-when-vectorising is a discipline.
\item \textbf{Memory-bound blindness.} At $n>2^{22}$ the kernel is
      bandwidth-limited; more lanes don't buy more bytes/sec.
\end{itemize}

\section{Summary}
\begin{tabularx}{\linewidth}{@{}l X@{}}
\toprule
\textbf{Piece} & \textbf{Verified version} \\
\midrule
lanes    & 4 floats per issue --- the cheap parallelism \\
saxpy    & scalar$\to$NEON = 3--4$\times$ real, BW-capped \\
auto-vec & compiler = intrinsics ($1.00\times$) at -O3 \\
SoA      & 2.35$\times$ for a pure layout change \\
RVV      & vector-length-agnostic --- the portable SIMD \\
\bottomrule
\end{tabularx}

\section*{Further reading}
\begin{itemize}
\item ARM NEON programmer's guide: \texttt{developer.arm.com} (SIMD ISAs)
\item RISC-V V spec (vector-length-agnostic): \texttt{riscv.org/specifications}
\item CS:APP ch5 (vectorisation within optimisation):
      \url{https://csapp.cs.cmu.edu/}
\end{itemize}

\begin{rulebox}[Watch instead]
\begin{itemize}
\item \textbf{SIMD intrinsics --- Computerphile} --- same
      experiment as this chapter on x86 (verified 2026-10-11).
      \url{https://youtu.be/x9Scb5Mku1g}
\end{itemize}
\end{rulebox}

\section*{Notes}
\begin{itemize}
\item All numbers verbatim from the chapter's two C programs on
      Apple Silicon; the $1.00\times$ auto-vec result is the
      chapter's second honest surprise after ch15's audit.
\item \texttt{rvv\_sum.S} from the outline is a read-along
      (RISC-V V encoding shown, not run --- no RVV hardware here).
\end{itemize}

\section*{Exercises}
\addcontentsline{toc}{section}{Exercises}
\stepcounter{section}

\begin{exercise}[easy]
Why is data-level parallelism ``free'' instruction-level
parallelism? One line.
\end{exercise}
\begin{solution}[hint]
One issue retires 4 lanes --- the pipeline never sees the other
three instructions.
\end{solution}
\begin{exercise}[easy]
Our scalar-vs-NEON gap is 3$\times$ at 4 lanes, not 4. What bounds
it at large n?
\end{exercise}
\begin{solution}[hint]
Memory bandwidth --- the roofline, ch9.
\end{solution}
\begin{exercise}[easy]
What does \texttt{-fno-vectorize} prove about our baseline?
\end{exercise}
\begin{solution}[hint]
That the ``scalar'' number is scalar, not the compiler's hidden
SIMD.
\end{solution}
\begin{exercise}[easy]
Which layout lets the vector unit fetch 4 useful floats per line:
AoS or SoA?
\end{exercise}
\begin{solution}[hint]
SoA --- AoS interleaves the fields you're not using.
\end{solution}
\begin{exercise}[medium]
Design the experiment that separates lane-parallel gain from
cache-friendly gain in the SoA result. What would you time?
\end{exercise}
\begin{solution}[hint]
Time AoS-with-2-fields (a stripped struct) vs SoA --- isolates
the vectorisation from the memory footprint.
\end{solution}
\begin{exercise}[medium]
Auto-vectorisation fails on which loop shapes? Name two and the
fix for each.
\end{exercise}
\begin{solution}[hint]
Data-dependent exits, and pointer aliasing (\texttt{restrict});
unrollable-count loops help.
\end{solution}
\begin{exercise}[medium]
Why does RVV's vector-length-agnostic loop beat NEON for portable
code? Trade-off?
\end{exercise}
\begin{solution}[hint]
Same source runs on 128/512/1024-bit hardware; NEON needs a
rewrite per width.
\end{solution}
\begin{exercise}[medium]
A kernel is 90\% memory-bound. Argue why intrinsics are the wrong
tool, and name the right one.
\end{exercise}
\begin{solution}[hint]
The bottleneck is bytes not lanes; fix the access pattern
(blocking, SoA, prefetch) instead.
\end{solution}
\begin{exercise}[hard]
Design the SoA+NEON+blocking version of a particle integrator and
predict the total speedup over the naive AoS scalar. What are the
two multiplicative terms?
\end{exercise}
\begin{solution}[hint]
Lane parallelism ($\approx$4$\times$) $\times$ footprint
($\approx$2$\times$) --- but they share the bandwidth ceiling;
predict sub-multiplicative.
\end{solution}
\begin{exercise}[hard]
The folklore entry said ``intrinsics beat compiler'' fails at
$1.00\times$. Write the counterexample where intrinsics DO win
(hint: ops the compiler can't express) and verify on your machine.
\end{exercise}
\begin{solution}[hint]
Shuffle/permute, saturating arithmetic, crypto primitives ---
the compiler won't find them from scalar code.
\end{solution}
\begin{exercise}[easy]
Rank the three parallelisms by cost-per-speedup for a
compute-bound kernel.
\end{exercise}
\begin{solution}[hint]
DLP (free, 4$\times$) $>$ ILP (already there) $>$ TLP (real
cost --- ch13's locks).
\end{solution}
\begin{exercise}[medium]
At what n does our saxpy become memory-bound? Estimate from the
table and the machine's bandwidth.
\end{exercise}
\begin{solution}[hint]
Gap narrows at $2^{25}$ --- about when the working set leaves
cache (ch9's mountain).
\end{solution}
\begin{exercise}[medium]
Explain why SoA is a vectorisation \emph{enabler}, not just a
layout.
\end{exercise}
\begin{solution}[hint]
Contiguous same-type arrays = clean lane loads; AoS forces
gathers or shuffles.
\end{solution}
\begin{exercise}[medium]
A teammate proposes AVX-512 ``for 4$\times$ more lanes''. Give
the two honest cautions beyond availability.
\end{exercise}
\begin{solution}[hint]
Frequency throttling on some parts, and bandwidth caps at large
n regardless of width.
\end{solution}
\begin{exercise}[hard]
Predict the RVV loop's behaviour at vlen=128 vs 1024 given our
bandwidth numbers --- at what vlen does the extra width stop
helping?
\end{exercise}
\begin{solution}[hint]
When the vector unit saturates memory BW --- on this class of
machine, above $\sim$256 bits.
\end{solution}
